\textbf{Barkad Achaari}, agent de police, subordonné de Lisnic. D'origine marocaine. Zélé, susceptible quand on parle de sa famille. Cousin de Moustapha Nedali.

\textbf{Tiago Alvarez}, travesti brésilien. Brûlé vif par les « Disciples de la Colère ». Sa mort déclenche l'enquête.

\textbf{Atiena} (Repentance Whittingham), « Gardienne de la Nuit ». Femme d'une beauté saisissante, d'origine khoïsan (Afrique australe). Officiellement femme de chambre au Caribbean Hôtel. Enlève Titus pour le protéger. Arrêtée avec 10 kg de cocaïne dans le coffre de sa Mini.

\textbf{Mathéo Baleya}, associé de Greg au cabinet de détectives privés.

\textbf{Rose Barbet} (née Lortie), 75 ans. Épouse soumise et épuisée. Découvre un colis contenant une tête humaine.

\textbf{Thierry Barbet}, mari de Rose. Alcoolique violent.

\textbf{Père Beaubois}, prêtre confesseur. Victime du Brain Catcher : tué, décapité. Son corps est découvert par Solange Jacquard.

\textbf{Titus Beaugendre}, lieutenant de police, meilleur ami et partenaire de Djeferson depuis l'enfance. Grand, noir, d'origine sierra-léonaise. Loyal et expéditif. Mari de Bérénice, père de Paul et Virginie. Enlevé par Atiena pour le protéger.

\textbf{Bérénice Beaugendre}, femme de Titus. Cuisinière remarquable, caractère bien trempé. Jalouse du lien fusionnel entre son mari et Djeferson. Soupçonnée par Djeferson d'une liaison avec Greg.

\textbf{Paul Beaugendre}, fils de Titus et Bérénice.

\textbf{Virginie Beaugendre}, fille de Titus et Bérénice. Studieuse, lit saint Augustin.

\textbf{Djeferson Beauvais}, narrateur et protagoniste. Commandant de police judiciaire. Cynique, amateur de cigares, de bons vins et de vieux livres. Règle volontiers les problèmes de façon expéditive. Armé de « Manu », son 6.35 Manufrance de collection hérité de son grand-père Philibert, résistant. Père militaire absent, prénom hérité d'un Thomas Jefferson mal orthographié.

\textbf{Gulav Behram} dit « le Kurde », criminel multirécidiviste : vol, viol, extorsion, cambriolage. Constamment libéré faute de place en prison. Trop violent même pour les codétenus. Tué par Djeferson lors d'un cambriolage.

\textbf{Ronny Bell}, ancien skinhead, tient le bar « Le Bouclier », repaire de l'extrême droite.

\textbf{Comte Léopold Chiasson de Bellisle} \emph{(décédé)}, aristocrate raffiné, homosexuel discret. Admirait Robert. Mort d'infarctus dans sa roseraie avant le début du récit.

\textbf{Alena Benesch}, infirmière, blonde aux yeux noisette. Victime du réseau de Charrier. A porté plainte puis s'est rétractée. Dans le coma à l'hôpital.

\textbf{Abraham Botrel}, 53 ans, « folle notoire ». Retrouvé mort dans une fosse à purin.

\textbf{Ernst Brestrich}, personnage fictif, auteur du « Traité sur les races ».

\textbf{Zarina Brizzi} (Fragale Di Brizzi), aristocrate italienne, pianiste de talent, excellente cuisinière. Brune, dort comme une masse et occupe tout le lit. Sœur jumelle de Tosca. Compagne de Djeferson.

\textbf{Tosca Brizzi}, sœur jumelle de Zarina (italienne). Blonde --- contrairement à Zarina, brune. Compagne de Zaahid Shirani. Présente dans l'intimité de l'appartement de Djeferson.

\textbf{Tacito Cerqueira}, associé de Jégou à l'entreprise de déménagement. Portugais corpulent, très poilu. Double vie : facho le jour, travesti prostitué la nuit. Parents tous deux transgenres. Force physique exceptionnelle. Lit \emph{Mein Kampf}. Devient informateur.

\textbf{Docteur Sébastien Charrier}, chirurgien, ancien champion universitaire de boxe. Propriétaire du château « La Frétoise ». Impliqué dans un réseau de pervers (affaire Alena Benesch). Ami de Rémi Durand.

\textbf{Professeur Wilfrid Chauveau}, neurochirurgien à la Pitié-Salpêtrière. A opéré Sam Girard après sa blessure par balle à la tête.

\textbf{Lou De La Croix} (Louise, « Loulou »), nouvelle compagne de Greg. Issue d'une vieille famille d'aristocrates désargentés ; père escroc notoire en fuite. Nymphomane, épuise Greg sexuellement. Absente des scènes d'action.

\textbf{Dardariel}, sbire et homme de main de Monseigneur Riqueti.

\textbf{Noé Desmarais}, membre des Disciples de la Colère. Sa mort est annoncée à Sally par Greg.

\textbf{Rémi Durand}, gynécologue. Retrouvé pendu dans son garage (mort suspecte). Ami et complice de Charrier.

\textbf{Monster Gang}, organisation fictive du roman.

\textbf{Samuel Girard} (Sam), capitaine, ancien des Forces Spéciales (Légion étrangère, missions en Syrie et Afghanistan). Blessé près de Kandahar --- balle dans la tête, opéré par le Pr Chauveau. Crises de tétanie passagères. Cuisinier hors pair. Ami de Djeferson.

\textbf{Père Marian Granet}, prêtre aux penchants BDSM. Capturé et torturé à la chaise électrique par Riqueti. Survit et est relâché sur intervention de Djeferson.

\textbf{Le réceptionniste du Caribbean Hôtel} (Dumo), métis du Cap (mère française, père khoïkhoï). Polyglotte, hypnotisant. Gardien d'un établissement rattaché à l'ambassade de Namibie, interdit aux Blancs. Protecteur d'Atiena.

\textbf{Solange Jacquard}, 90 ans. Dévote acharnée, homosexuelle réprimée. Découvre le corps du père Beaubois.

\textbf{Marc-Antoine Jacquinot}, autre voisin au 6e étage, « l'ermite du sixième ». Professeur de philosophie, solitaire, collectionneur de livres anciens rares (certains lui seront volés). Refuse toutes les corvées de voisinage.

\textbf{Aymeric Jégou}, chef du groupe néonazi « Les Disciples de la Colère ». Entrepreneur en déménagement. Blond aux yeux bleus, sec, nerveux et vicieux. Responsable du meurtre de Tiago Alvarez.

\textbf{Simon Keskula}, personnage fictif, lié à l'Alliance de la Révélation.

\textbf{Reckless \& Knot}, société financière fictive. Ancien employeur de Greg Lussier (analyste M\&A).

\textbf{Korax}, chat roux tigré de la mère Ouvrard. Sournois et territorial, sait utiliser l'ascenseur. Déteste Djeferson.

\textbf{Brigadier-chef Darian Lisnic}, policier en uniforme, d'origine roumaine. Borné, raciste, soupçonneux. Chef du duo qui intercepte le groupe de Djeferson.

\textbf{Grégoire Lussier} (Greg), détective privé, ancien analyste M\&A chez Reckless \& Knot. Frère aîné de Nathan. Bavard, intellectuel de façade. Ami et auxiliaire de Djeferson. Liaison probable avec Bérénice.

\textbf{Nathan Lussier}, frère de Greg. Garagiste autodidacte, modifie la Kangoo de Djeferson (250 cv, blindage, humidor dans la boîte à gants). Passionné de bel canto, de cinéma et de chasse à l'arc. Pilote le van G30 lors des expéditions.

\textbf{Maître Yoshimura Masayoshi}, personnage fictif, à Kobe.

\textbf{Louis-Marie de la Maîtraie}, éleveur de chevaux, particule nobiliaire. Pédophile et cannibale, auteur de massacres d'enfants du voisinage. Tabassé par Djeferson et Titus.

\textbf{Jessica Millet}, secrétaire de l'entreprise de déménagement de Jégou.

\textbf{Milo Monteil}, membre des Disciples de la Colère.

\textbf{Niccolo}, garde du corps de Riqueti. Gigantesque, oreilles en chou-fleur, ancien boxeur. Exécutant violent et fidèle. Tué par Titus.

\textbf{Ordre de la Lune Noire}, organisation fictive du roman.

\textbf{Maria Ouvrard} (« mère Ouvrard »), voisine de palier de Djeferson au 6e étage. Maîtresse du chat Korax.

\textbf{Dylan Passereau}, chauffeur routier colossal (plus de 2 m). Peu dégourdi. Reçoit un colis contenant des organes génitaux ; suspect provisoire.

\textbf{Jean-Charles Pichard} (domaine Meursault Chevalières), vigneron fictif.

\textbf{Robert Pleimelding}, garde-chasse et propriétaire terrien. Solitaire, chasseur passionné. Époux de Claire, maître de la chienne Greta (drahthaar). Découvre un cadavre calciné en forêt. Héritier du comte de Bellisle.

\textbf{Claire Pleimelding}, épouse de Robert. Cuisinière d'exception. Mariage fondé sur la discrétion et la confiance mutuelle.

\textbf{Romuald Pueyrredon}, associé de Greg au cabinet de détectives privés.

\textbf{Valérie Renou}, retraitée. Découvre l'un des colis suspects et meurt d'une crise cardiaque.

\textbf{Monseigneur Mathéo Riqueti}, cardinal, évêque. Beau et raffiné en apparence. Véritable identité du « Brain Catcher » : tue et décapite des prêtres pédophiles, emplit leurs crânes de croquettes de luxe --- vengeance contre l'Église qui avait fait tuer son chien d'enfance.

\textbf{Sally Robinson}, meneuse de revue au cabaret Sugar \& Spice. Femme trans. Cliente de Greg, participe aux expéditions nocturnes.

\textbf{Roman}, médecin (prénom seulement). Viole ses patientes sous couvert d'examens gynécologiques. Mentionné par Charrier.

\textbf{Le Narcisse Rose}, bar/club fictif du roman.

\textbf{L'Alliance de la Révélation}, communauté fictive.

\textbf{Matthias Schuster}, curé de Montaulogne, personnage fictif.

\textbf{Docteur Zaahid Shirani}, médecin légiste, ami et confident de Djeferson. D'origine bengalie. Fumeur de cannabis. Expert des autopsies. Compagnon de Tosca Brizzi.

\textbf{Terzo}, cirneco de l'Etna, chien de Riqueti. Lévrier sicilien de haute lignée. Tué par Djeferson.

\textbf{Père Clément Vidal}, curé d'une quarantaine d'années. Pédophile. Première victime du « Brain Catcher » : disparaît et est décapité.
